\section{Plug-and-play methods and Regularization by denoising}

\subsection{MAP cost function and proximal algorithms}

Let use consider the linear inverse problem from \eqref{eq:inverse-problem}
and let $p(y|x)$ be the likelihood of the measurements $y$ given the unknown image $x$, and $p(x)$ be the prior distribution of the unknown image.
The goal is to find an estimate $\hat{x}$ of the unknown image $x$ that maximizes the posterior distribution $p(x|y)$, which can be expressed using Bayes' theorem as:
\begin{equation*}
    \hat{x} = \arg \max_x p(x|y) = \arg \max_x p(y|x)p(x).
\end{equation*}
Taking the negative logarithm, we obtain the following minimization problem:
\begin{equation}
    \hat{x} = \arg \min_x f(x, y) - \beta s(x)
    \label{eq:map}
\end{equation}
where $-\ln p(x) = s(x) + C$ and under the likelihood  $p(y | x) \propto \exp(f(x, y))$, where $f(x)$ is the data fidelity term and $\beta$ is a hyperparameter balancing the two terms.
The choice of the function $s$ is crucial as it encodes the prior knowledge about the unknown image $x$ and thus defines the regularization of the inverse problem.
One of the most common choices for $s$ is the Total Variation (TV) regularization, which promotes piecewise constant images and is defined as the $\ell_1$ norm of the image gradient \cite{Rudin1992}. Another popular choice is the Tikhonov regularization, which promotes smooth images and is defined as the squared $\ell_2$ norm of a linear transformation of the image \cite{tikhonov1977solutions}.

\subsubsection{Proximal algorithms}

\begin{definition}
    The proximal operator of a function $g$ is defined as:
    \begin{equation*}
        \operatorname{prox}_g(z) = \arg \min_x \frac{1}{2} \|x - z\|_2^2 + g(x).
    \end{equation*}
\end{definition}

Generally, the solution of the MAP problem in \eqref{eq:map} does not have a closed-form expression and must be found using iterative optimization algorithms.
A common strategy is to split the optimization problem into two sub-problems, one for the data fidelity term and one for the regularization term, and iteratively solve them while enforcing consistency between the two solutions.
\cite{https://doi.org/10.48550/arxiv.0912.3522} propose a first order splitting algorithm to solve the proximal problem, which consists in performing a gradient descent step on the data fidelity term followed by a proximal step on the regularization term.
We can also mention the Half Quadratic Splitting (HQS) algorithm \cite{392335}.
\begin{algorithm}
    \caption{Half Quadratic Splitting (HQS)}
    \label{alg:hqs}
    \begin{algorithmic}[1]
        \Require initial $x_0$, measurements $y$, penalty parameter $\rho>0$, max iterations $K$
        \For{$k=0,\dots,K-1$}
            \State $z_k \gets \arg \min_z f(z,y) + \frac{\rho}{2} \|z - x_k\|_2^2$
            \State $x_{k+1} \gets \arg \min_x s(x) + \frac{\rho}{2} \|x - z_k\|_2^2$
        \EndFor
        \Ensure $x_K$
    \end{algorithmic}
\end{algorithm}
Then the Alternating Direction Method of Multipliers (ADMM) \cite{boyd_distributed_2010} is a more general algorithm that can be applied to a wider range of problems, including non-convex problems, and has been shown to have good convergence properties.
The idea is to use the augmented Lagrangian to enforce the consistency between the two sub-problems.
% \begin{algorithm}
%     \caption{Plug-and-Play Forward--Backward Splitting (PnP-FBS)}
%     \label{alg:pnp-fbs}
%     \begin{algorithmic}[1]
%         \Require initial $x_0$, measurements $y$, step size $\tau>0$, regularization weight $\lambda$, max iterations $K$
%         \For{$k=0,\dots,K-1$}
%             \State $z_k \gets x_k - \tau \nabla_x f(x_k,y)$
%             \State $x_{k+1} \gets \operatorname{prox}_{\tau\lambda s}(z_k)$ \Comment{or replace prox by a denoiser $D_{\sigma}$ for PnP}
%         \EndFor
%         \Ensure $x_K$
%     \end{algorithmic}
% \end{algorithm}
% ADMM algorithm
\begin{algorithm}
    \caption{ADMM}
    \label{alg:pnp-admm}
    \begin{algorithmic}[1]
        \Require initial $x_0$, $v_0$, $u_0$, measurements $y$, penalty parameter $\rho>0$, max iterations $K$
        \For{$k=0,\dots,K-1$}
            \State $x_{k+1} \gets \arg \min_x f(x,y) + \frac{\rho}{2} \|x - v_k + u_k\|_2^2$
            \State $v_{k+1} \gets \operatorname{prox}_{\frac{1}{\rho} s}(x_{k+1} + u_k)$ %\Comment{or replace prox by a denoiser $D_{\sigma}$ for PnP}
            \State $u_{k+1} \gets u_k + x_{k+1} - v_{k+1}$
        \EndFor
        \Ensure $x_K$
    \end{algorithmic}
\end{algorithm}

\subsection{Plug-and-Play methods}

\subsection{Regularization by denoising}

Regularization by denoising (RED) is a framework that leverages the power of denoising algorithms to solve problems of the form \eqref{eq:map} \cite{RED_romano}.
The core idea is to replace $s$ with an explicit regularization function defined as:
\begin{equation*}
    s(x) = \frac{1}{2} x^T (x - D(x))
\end{equation*}
where $D$ is a denoising operator.
The motivation behind this choice is that the regularization term encourages the solution to be close to its denoised counterpart.

\cite{RED_romano} show that image denoisers are locally homogeneous (LH), \textit{i.e.} $D((1+ \alpha) x) = (1+ \alpha) D(x)$ for sufficiently small $\alpha \neq 0$.
Under this condition, they claim that the regularizer obey the gradient rule :
\begin{equation*}
    \nabla s(x) = x - D(x).
\end{equation*}
and propose several iterative algorithms to solve the resulting optimization problem under quadratic loss as the data-fidelity term, including gradient descent, ADMM and the fixed-point strategy.

\cite{RED_clarifications} later show that the LH condition is not sufficient to guarantee the gradient rule, and that the denoiser must also have a symmetric Jacobian matrix, \textit{i.e.} $J_D(x) = J_D(x)^T$ for all $x$.
This condition is not satisfied by many popular denoisers, such as BM3D \cite{BM3D} and DnCNN \cite{DnCNN}, which limits the applicability of the RED framework.